\documentclass[11pt,letterpaper]{article}
\usepackage[T1]{fontenc}
\usepackage{lmodern}
\usepackage{amsmath,amssymb,amsthm}
\usepackage[margin=1.08in,nohead,includefoot]{geometry}
\usepackage[hidelinks,pdfencoding=auto]{hyperref}
\hypersetup{pdftitle={Weak separation and finite-trace coding: independence of infinite fin-intersecting MAD families},pdfauthor={Haoxuan Ye}}
\newtheorem{theorem}{Theorem}[section]
\newtheorem{lemma}[theorem]{Lemma}
\newtheorem{corollary}[theorem]{Corollary}
\newtheorem{proposition}[theorem]{Proposition}
\theoremstyle{definition}
\newtheorem{definition}[theorem]{Definition}
\theoremstyle{remark}
\newtheorem{remark}[theorem]{Remark}
\newcommand{\s}{\mathfrak{s}}
\newcommand{\bnum}{\mathfrak{b}}
\newcommand{\anum}{\mathfrak{a}}
\newcommand{\ap}{\mathfrak{ap}}
\newcommand{\dpnum}{\mathfrak{dp}}
\newcommand{\cont}{\mathfrak{c}}
\newcommand{\A}{\mathcal{A}}
\newcommand{\M}{\mathcal{M}}
\newcommand{\U}{\mathcal{U}}
\newcommand{\FI}{\operatorname{FI}}
\newcommand{\Con}{\operatorname{Con}}
\newcommand{\tr}{\operatorname{tr}}
\newcommand{\Id}{\mathcal{I}}
\newcommand{\E}{\mathsf{E}}
\newcommand{\PP}{\mathsf{P}}
\setlength{\emergencystretch}{1em}
\title{Weak separation and finite-trace coding:\\ independence of infinite fin-intersecting MAD families}
\author{Haoxuan Ye}
\date{29 September 2026}
\begin{document}
\maketitle
\begin{abstract}
We prove that, if $\s<\ap$, an almost disjoint family is fin-intersecting
exactly when its cardinality is less than $\s$. Consequently, no infinite
fin-intersecting maximal almost disjoint family exists in this region.
The main lemma converts countably many weak intersection tests on a family
of size less than $\bnum$ into eventual traces of one sequence of disjoint
nonempty finite sets. Weak separation then encodes a splitting family and
its complements on prescribed members of any sufficiently large almost
disjoint family. Combined with a consistency theorem of Banakh, Machura
and Zdomskyy and the positive results of Corral and Rodrigues, this yields
the relative independence over ZFC of the existence of an infinite
fin-intersecting MAD family. The negative model also contains a MAD family
with pseudocompact Vietoris hyperspace. A companion Lean development
checks the exact obstruction, the comparison $\ap\leq\bnum$, and the CH
existence construction; the forcing models and the full ZFC independence theorem
are not formalized there.
\end{abstract}

\section{Introduction}
An almost disjoint family can have centered traces along an infinite
subsequence of every sequence of disjoint finite blocks. Corral and
Rodrigues call this property \emph{fin-intersecting} and use it to obtain
MAD families whose associated Vietoris hyperspaces are pseudocompact
\cite{CR}. Their Question~3.7 asks whether ZFC proves the existence of a
fin-intersecting MAD family. Throughout this paper, a MAD family is
required to have infinitely many members.

Let $\E$ denote the assertion that an infinite fin-intersecting MAD family
on $\omega$ exists. The main combinatorial result is the following.
\begin{theorem}\label{thm:main}
Assume $\s<\ap$. For every almost disjoint family $\A\subseteq[\omega]^\omega$,
\[
 \FI(\A)\quad\Longleftrightarrow\quad |\A|<\s.
\]
In particular, $\neg\E$ holds.
\end{theorem}

The positive direction of the equivalence is the known small-family
bound of Corral and Rodrigues. Our argument for the converse starts with
an arbitrary candidate family. We choose pairs of distinct members indexed
by a splitting family and prescribe complementary traces on each pair.
Weak separation supplies countably many infinite testing sets. A single
eventual bound converts these tests into disjoint finite blocks. Any
infinite set of indices is split by one of the labels, giving two infinite
retained traces whose intersection is finite.

The inequality $\ap\leq\bnum$ is the external cardinal comparison that makes
this coding possible. Banakh, Machura and Zdomskyy give a relative
consistency model with $\s_\omega=\aleph_1<\dpnum=\cont=\aleph_2$ \cite{BMZ}.
It satisfies $\s<\ap$, so Theorem~\ref{thm:main} rules out every infinite
FI MAD family there. On the other hand, Corral and Rodrigues prove
$\ap=\s=\cont\Rightarrow\E$, in particular $\mathrm{CH}\Rightarrow\E$.
Thus, with consistency understood syntactically,
\[
 \Con(\mathrm{ZFC})\ \Longrightarrow\
 \Con(\mathrm{ZFC}+\E)\ \wedge\ \Con(\mathrm{ZFC}+\neg\E).
\]

Sections~\ref{sec:coding}--\ref{sec:obstruction} contain the new coding and
exclusion arguments. Section~\ref{sec:ch} gives a self-contained reconstruction
of the known CH construction. Section~\ref{sec:independence} states the model
inputs precisely and derives independence. We also separate FI existence
from the existence of a MAD family with pseudocompact hyperspace.
Section~\ref{sec:formal} records exactly what the companion Lean files prove.

\section{Definitions and preliminary facts}
We work in ZFC. Write $A\subseteq^*B$ when $A\setminus B$ is finite and
$A=^*B$ when their symmetric difference is finite. For functions on $\omega$,
$f\leq^*g$ means that $f(n)\leq g(n)$ for all sufficiently large $n$.
We use $[0,k)=\{0,\ldots,k-1\}$.

\begin{definition}
A family $\A\subseteq[\omega]^\omega$ is \emph{almost disjoint} (AD) if distinct
members have finite intersection. It is \emph{MAD} if $\A$ is infinite and
\[
 \forall X\in[\omega]^\omega\ \exists A\in\A\quad |X\cap A|=\aleph_0.
\]
Finite or empty AD families are allowed in intermediate arguments, but do
not count as MAD families.
A \emph{fin-sequence} is a sequence $C=(C_n:n<\omega)$ of pairwise disjoint
nonempty finite subsets of $\omega$. Its trace on $A$ is
\[
 \tr_C(A)=\{n:A\cap C_n\ne\varnothing\}.
\]
There is no uniform bound on the block sizes, and the blocks need not
cover $\omega$.
\end{definition}

\begin{definition}[{\cite[Definition~2.3]{CR}}]
A family of infinite sets is \emph{centered} if every nonempty finite
subfamily has infinite intersection. The empty family is centered.
An AD family $\A$ is \emph{fin-intersecting} (FI) if, for every fin-sequence
$C$, some infinite $I\subseteq\omega$ makes
\[
 \{I\cap\tr_C(A):A\in\A,\ |I\cap\tr_C(A)|=\aleph_0\}
\]
centered. The test applies to \emph{all} infinite retained traces, after
finite traces have been discarded.
\end{definition}

The splitting number $\s$ is the least size of a family $\mathcal S$ of
subsets of $\omega$ such that every infinite $I$ is split by some
$S\in\mathcal S$: both $I\cap S$ and $I\setminus S$ are infinite.
The bounding number $\bnum$ is the least size of an unbounded family in
$(\omega^\omega,\leq^*)$. Let $\anum$ be the least size of an infinite MAD
family, and let $\cont=2^{\aleph_0}$.

For $\mathcal B\subseteq\A$, a set $X\subseteq\omega$ \emph{weakly separates}
$\mathcal B$ from $\A\setminus\mathcal B$ if
\[
 |A\cap X|=\aleph_0\ (A\in\mathcal B),\qquad
 |A\cap X|<\aleph_0\ (A\in\A\setminus\mathcal B).
\]
The invariant $\ap$ is the least size of an AD family with a subfamily
that cannot be weakly separated from its complement. This definition uses
infinite versus finite intersection, not almost containment.

We use the standard inequalities
\begin{equation}\label{eq:bounds}
 \aleph_1\leq\ap\leq\bnum\leq\anum\leq\cont,
 \qquad \aleph_1\leq\s\leq\cont.
\end{equation}
In particular, $\ap\leq\bnum$ follows from \cite[Theorem~2]{BMZ}.
For completeness we recall the diagonal argument for $\bnum\leq\anum$.

\begin{lemma}\label{lem:ba}
An infinite AD family of size less than $\bnum$ is not MAD.
\end{lemma}
\begin{proof}
Choose distinct $A_i\in\A$ and put $E_i=A_i\setminus\bigcup_{j<i}A_j$.
The sets $E_i$ are infinite and pairwise disjoint. For each $A\in\A$, let
\[
 f_A(i)=
 \begin{cases}
 1+\max(\{0\}\cup(A\cap E_i)),&A\ne A_i,\\
 0,&A=A_i.
 \end{cases}
\]
Choose a common eventual bound $g$ and choose $x_i\in E_i$ with $x_i\geq g(i)$.
The points are distinct. For a fixed $A$, all sufficiently late $x_i$ lie
outside $A$, except possibly the unique index with $A=A_i$.
Thus $\{x_i:i<\omega\}$ is infinite and almost disjoint from every member.
\end{proof}

\begin{lemma}\label{lem:small}
The FI property passes to subfamilies. Every AD family of size less than
$\s$ is FI.
\end{lemma}
\begin{proof}
For heredity, every retained trace for a subfamily is a retained trace for
the original family. For the second assertion, fix $C$. Its trace family
has size less than $\s$, so it fails to split some infinite $I$.
Each trace is finite or cofinite in $I$. Every retained trace is therefore
cofinite in $I$, and the retained family is centered.
This is the argument of \cite[Theorem~3.1]{CR}.
\end{proof}

\section{Uniform finite-trace coding}\label{sec:coding}
The following lemma makes no almost-disjointness assumption. Its point is
to use one eventual bound before choosing any of the block endpoints.

\begin{lemma}[Uniform coding]\label{lem:coding}
Let $\U\subseteq\mathcal P(\omega)$ have size less than $\bnum$, and assign
$S_A\subseteq\omega$ to each $A\in\U$. Suppose that infinite sets
$X_n\subseteq\omega$ satisfy
\[
 |A\cap X_n|=\aleph_0\quad\Longleftrightarrow\quad n\in S_A.
\]
There is a fin-sequence $C$ such that $\tr_C(A)=^*S_A$ for every $A\in\U$.
\end{lemma}
\begin{proof}
First bound each negative test. Set
\[
 d_A(k)=
 \begin{cases}
 0,&k\in S_A,\\
 1+\max(\{0\}\cup(A\cap X_k)),&k\notin S_A.
 \end{cases}
\]
Thus, in the negative case, $A\cap X_k\subseteq[0,d_A(k))$.
For a positive test and an endpoint $k$, the infinite set $A\cap X_i$ has a
least element at least $k$. Collect the finitely many relevant witnesses:
\[
 P_A(k)=\{\min(A\cap X_i\cap[k,\infty)):i\leq k,\ i\in S_A\},
 \quad h_A(k)=\max(\{k+1,d_A(k)\}\cup P_A(k)).
\]
Choose $g\in\omega^\omega$ with $h_A\leq^*g$ for every $A\in\U$, increasing
it if necessary so that $g(k)\geq k+1$. This is the sole use of $|\U|<\bnum$.

Define the endpoints recursively. Let $r_0=0$ and, given $r_n$, put
\begin{align*}
 \ell_n&=\max\{n,g(n),r_n\},\\
 z_n&=\min(X_n\cap[\ell_n,\infty)),\\
 r_{n+1}&=1+\max\{\ell_n,g(\ell_n),z_n\},
 &C_n&=X_n\cap[\ell_n,r_{n+1}).
\end{align*}
Every $z_n$ exists and belongs to $C_n$, so the blocks are nonempty and finite.
Since $\ell_{n+1}\geq r_{n+1}$, the successive windows, and hence all the
blocks, are disjoint.

Fix $A$ and choose $N_A$ so that $h_A(k)\leq g(k)$ for every $k\geq N_A$.
For $n\geq N_A$ with $n\notin S_A$, we have
\[
 A\cap X_n\subseteq[0,d_A(n))\subseteq[0,g(n))\subseteq[0,\ell_n),
\]
so $A\cap C_n=\varnothing$. If $n\in S_A$, then $n\leq\ell_n$ and
$\ell_n\geq N_A$. The witness for test $n$ occurs in $P_A(\ell_n)$, giving
\[
 \ell_n\leq\min(A\cap X_n\cap[\ell_n,\infty))
 \leq h_A(\ell_n)\leq g(\ell_n)<r_{n+1}.
\]
It lies in $A\cap C_n$. Thus membership in $\tr_C(A)$ and $S_A$ agrees from
$N_A$ onward, as required.
\end{proof}

\begin{remark}
All functions $h_A$ are defined before $g$ and the endpoints are chosen.
The later evaluation at $\ell_n$ is legitimate because $\ell_n\geq n$.
The thresholds $N_A$ can depend on $A$; no uniform threshold over the whole
family is asserted. Any choices of positive witnesses in place of the
minima also work, which is the formulation used in Lean.
\end{remark}

\begin{corollary}[Arbitrary labels on a small infinite AD family]\label{cor:labels}
If $\A$ is infinite, AD, and $|\A|<\ap$, then every assignment
$A\mapsto S_A\subseteq\omega$ is realized modulo finite by the traces of
one fin-sequence.
\end{corollary}
\begin{proof}
For each $n$, put $\mathcal B_n=\{A\in\A:n\in S_A\}$.
If $\mathcal B_n$ is nonempty, weak separation provides a testing set $X_n$;
its infinite intersection with any positive member makes $X_n$ infinite.
If $\mathcal B_n$ is empty, Lemma~\ref{lem:ba} and $|\A|<\ap\leq\bnum$
provide an infinite set almost disjoint from all of $\A$; use it as $X_n$.
Apply Lemma~\ref{lem:coding}.
\end{proof}

Infinitude of the family in this corollary matters: a finite partition of
$\omega$ cannot have all its traces eventually empty. In the main argument
below every positive fiber is nonempty, so this special case can also be avoided.

\section{Excluding every sufficiently large candidate}\label{sec:obstruction}
\begin{proof}[Proof of Theorem~\ref{thm:main}]
The forward obstruction is the only new direction. Let $\M$ be any AD
family with $|\M|\geq\s$. Since $2\s=\s$, choose distinct members
\[
 \A=\{a_\xi^0,a_\xi^1:\xi<\s\}\subseteq\M.
\]
Take a splitting family $(S_\xi:\xi<\s)$ and label the selected members by
$S_{a_\xi^0}=S_\xi$ and $S_{a_\xi^1}=\omega\setminus S_\xi$.
Corollary~\ref{cor:labels} gives a fixed fin-sequence $C$ with these traces
modulo finite. Equivalently, one can weakly separate, for each $n$, the
nonempty subfamily
\[
 \{a_\xi^0:n\in S_\xi\}\ \cup\ \{a_\xi^1:n\notin S_\xi\}
\]
and apply Lemma~\ref{lem:coding} directly using $\s<\ap\leq\bnum$.

Let $I$ be any infinite set of indices, and choose $S_\xi$ splitting it.
The finite errors
\[
 D_0=\tr_C(a_\xi^0)\mathbin\triangle S_\xi,\qquad
 D_1=\tr_C(a_\xi^1)\mathbin\triangle(\omega\setminus S_\xi)
\]
do not destroy infinitude of either restricted trace. Both therefore
survive the deletion of finite traces. But
\[
 \tr_C(a_\xi^0)\cap\tr_C(a_\xi^1)\subseteq D_0\cup D_1,
\]
so their intersection is finite. They contradict centeredness, and $I$
cannot witness FI for $C$. The same $C$ works for every infinite $I$.
Thus $\M$ is not FI. Lemma~\ref{lem:small} supplies the opposite implication.

Finally, any infinite MAD family has size at least $\anum\geq\bnum\geq\ap>\s$,
so it is excluded by the same argument.
\end{proof}

\begin{corollary}\label{cor:necessary}
ZFC proves $\E\Rightarrow\ap\leq\s$.
\end{corollary}
The quantifiers in the negative conclusion are $\forall\M\,\exists C\,\forall I$:
the blocking sequence can depend on the candidate family. This neither
changes its members nor merely constructs one special non-FI MAD family.

\section{The known positive direction under CH}\label{sec:ch}
Corral and Rodrigues prove the stronger positive theorem under
$\ap=\s=\cont$ \cite[Theorem~3.4]{CR}. We reconstruct the CH case to make
the two sides of the independence argument explicit; no novelty is claimed
for this positive direction.

For an AD family $\A$, let
\[
 \Id(\A)=\{Y\subseteq\omega:\exists F\in[\A]^{<\omega},\ Y\subseteq^*\bigcup F\}.
\]
Call $Y$ \emph{positive} if $Y\notin\Id(\A)$.

\begin{lemma}\label{lem:hit}
If $\A$ is countably infinite and AD, and $\mathcal Y$ is a countable family
of $\Id(\A)$-positive sets, there is an infinite set $a$ almost disjoint
from every member of $\A$ and meeting every $Y\in\mathcal Y$ infinitely.
\end{lemma}
\begin{proof}
Add $\omega$ to the requirements. It is positive: a finite subfamily cannot
almost cover an infinite AD family. Enumerate $\A=(A_i:i<\omega)$ and the
requirements as $(Y_j:j<\omega)$, repeating entries when necessary.
At stage $n$, for every $j\leq n$, choose a fresh point
$x_{j,n}\in Y_j\setminus\bigcup_{i\leq n}A_i$, avoiding all previously
chosen points. Positivity ensures infinitely many choices at each step.
The set of all these points meets each requirement infinitely and each
$A_i$ only in the finitely many points selected before stage $i$.
\end{proof}

\begin{lemma}\label{lem:cover}
Let $\A$ be AD and $C$ a fin-sequence. If an infinite $I$ satisfies
$\bigcup_{n\in I}C_n\in\Id(\A)$, some infinite $J\subseteq I$ witnesses FI
for $\A$ and $C$.
\end{lemma}
\begin{proof}
Choose finite $F\subseteq\A$ almost covering that union. Remove the finitely
many blocks meeting the finite uncovered part. Every remaining block is
contained in $\bigcup F$. There are only finitely many intersection
patterns with $F$, so one nonempty pattern $K\subseteq F$ occurs on an
infinite set $J$ of blocks. Members of $K$ have trace $J$ after restriction;
members of $F\setminus K$ have empty trace. Each $A\in\A\setminus F$ meets
$\bigcup F$ finitely, hence meets only finitely many of the disjoint blocks.
The only retained trace is therefore $J$, and centeredness follows.
This is the mechanism of \cite[Lemma~2.4]{CR}.
\end{proof}

\begin{theorem}[Known CH construction]\label{thm:ch}
Under CH, there is an infinite FI MAD family of size $\aleph_1$.
\end{theorem}
\begin{proof}
Start with a countably infinite AD family $(a_n:n<\omega)$. Using CH,
enumerate all infinite subsets as $(X_\alpha:\omega\leq\alpha<\omega_1)$
and all fin-sequences as $(C^\alpha:\omega\leq\alpha<\omega_1)$.
At stage $\alpha$, the preceding family
$\A_\alpha=\{a_\xi:\xi<\alpha\}$ is countably infinite.
Choose an infinite $I_\alpha$ on which every trace of an old member against
$C^\alpha$ is finite or cofinite. Successive thinning followed by a
diagonal pseudointersection gives such an $I_\alpha$.

For each already registered $\omega\leq\gamma\leq\alpha$ and finite
$F\subseteq\alpha$, put
\[
 S_\gamma(F)=I_\gamma\cap\bigcap_{\xi\in F}\tr_{C^\gamma}(a_\xi),
 \qquad Y_\gamma(F)=\bigcup_{n\in S_\gamma(F)}C_n^\gamma.
\]
Include in the requirement list every such $Y_\gamma(F)$ that is currently
$\Id(\A_\alpha)$-positive, together with $\omega$ and, if positive, $X_\alpha$.
This list is countable. Lemma~\ref{lem:hit} gives a new member $a_\alpha$
meeting all requirements infinitely and almost disjoint from the preceding
family. In particular,
\begin{equation}\label{eq:invariant}
 Y_\gamma(F)\notin\Id(\A_\alpha)
 \quad\Longrightarrow\quad |S_\gamma(F\cup\{\alpha\})|=\aleph_0.
\end{equation}
The implication uses the finiteness and disjointness of the blocks.

The resulting $\M=\{a_\xi:\xi<\omega_1\}$ is AD with distinct infinite
members, hence of size $\aleph_1$. It is maximal: at the stage for an
infinite $X_\alpha$, either finitely many preceding members almost cover it,
so one already meets it infinitely, or the new member is required to do so.

Fix $C=C^\gamma$. If any infinite union of its blocks belongs to $\Id(\M)$,
Lemma~\ref{lem:cover} gives an FI witness. Otherwise every infinite block
union is positive for the final ideal and consequently for every earlier
ideal. We show that $I_\gamma$ itself is a witness.
Take finitely many members with infinite traces on $I_\gamma$.
Those indexed below $\gamma$ have cofinite traces there, so their finite
intersection $L$ is infinite; if there are none, take $L=I_\gamma$.
Add the remaining indices in increasing order. Whenever the current
intersection $L$ is infinite, its block union is positive in the final
ideal, hence positive when the next member was added.
Condition~\eqref{eq:invariant} preserves infinitude of the intersection.
After finitely many steps the intersection of the chosen traces is infinite.
Thus all retained traces are centered and $\M$ is FI.
\end{proof}

\section{Relative independence and the external model}\label{sec:independence}
The invariant $\dpnum$ is the least size of the union of two orthogonal
families of infinite subsets of $\omega$ that cannot be weakly separated.
Orthogonality requires finite intersections between the two families.
An inseparable cut in an AD family is such a pair, so $\dpnum\leq\ap$.
Let $\s_\omega$ be the least size of a family containing a simultaneous
splitter for every countable sequence of infinite subsets of $\omega$.
In particular, $\s\leq\s_\omega$.

\begin{theorem}[External consistency input, {\cite[Theorem~4(3)]{BMZ}}]\label{thm:bmz}
Relative to the consistency of ZFC, it is consistent that
$\s_\omega=\aleph_1$ and $\dpnum=\cont=\aleph_2$.
\end{theorem}
This is a weakening of the configuration stated in the cited theorem.
Its proof refers to Dow's model and establishes the value of $\s_\omega$.
We use that published consistency result and do not reconstruct its forcing.
In this model, the inequalities give
\[
 \s=\aleph_1,\qquad \ap=\bnum=\anum=\cont=\aleph_2.
\]
Thus Theorem~\ref{thm:main} applies internally to every infinite MAD family.

\begin{theorem}\label{thm:independence}
If ZFC is consistent, then both $\mathrm{ZFC}+\E$ and
$\mathrm{ZFC}+\neg\E$ are consistent. Consequently $\E$ is independent of ZFC.
\end{theorem}
\begin{proof}
For the positive side use the standard relative consistency of CH with
ZFC and Theorem~\ref{thm:ch}. For the negative side use
Theorem~\ref{thm:bmz} and the ZFC implication in Theorem~\ref{thm:main}.
If either $\E$ or its negation were a theorem of ZFC, the respective
opposite extension would be inconsistent.
\end{proof}
This is a relative-consistency argument in the usual syntactic sense.
It does not assert that ZFC proves its own consistency, and does not assume
an extra transitive-model hypothesis. It answers the published existence
question \cite[Question~3.7]{CR}; in the earlier arXiv v2 this is Question~17.

\section{Consequences and limits}
\begin{corollary}\label{cor:section}
ZFC proves $\ap=\cont\Rightarrow(\E\leftrightarrow\s=\cont)$.
\end{corollary}
\begin{proof}
If $\s<\cont=\ap$, Theorem~\ref{thm:main} gives $\neg\E$.
If $\s=\cont=\ap$, apply \cite[Theorem~3.4]{CR}.
\end{proof}

For an AD family $\A$, the space $\Psi(\A)$ has underlying disjoint union
$\omega\mathbin{\dot\cup}\A$. Natural numbers are isolated; a basic
neighborhood of $A\in\A$ is $\{A\}\cup(A\setminus F)$ for finite $F$.
Write $\exp(X)$ for the hyperspace of nonempty closed subsets with the
Vietoris topology, generated by $U^+=\{F:F\subseteq U\}$ and
$U^-=\{F:F\cap U\ne\varnothing\}$ for open $U$.
Pseudocompactness means that every continuous real-valued function is bounded.
Let $\PP$ assert that some infinite MAD family $\A$ has pseudocompact
$\exp(\Psi(\A))$.

Corral and Rodrigues prove $\E\Rightarrow\PP$ \cite[Proposition~2.5]{CR}.
They also prove $\ap=\cont\Rightarrow\PP$ \cite[Theorem~5]{CRnew}.
\begin{corollary}\label{cor:topology}
If ZFC is consistent, then $\mathrm{ZFC}+\PP+\neg\E$ is consistent.
\end{corollary}
\begin{proof}
The model of Theorem~\ref{thm:bmz} satisfies $\ap=\cont$ and $\s<\ap$.
Apply the two preceding implications.
\end{proof}
This is stronger than producing one pseudocompact MAD family that is not FI:
there is no infinite FI MAD family in that model at all. It does not prove
independence of $\PP$, nor produce a model of $\neg\PP$.

The known sufficient condition $\anum<\s$ and our necessary condition give
\[
 \anum<\s\quad\Longrightarrow\quad\E\quad\Longrightarrow\quad\ap\leq\s.
\]
We do not establish sufficiency of $\ap\leq\s$, or even of $\ap=\s$.
The negative model refutes the implication from $\ap=\cont$ alone to FI MAD
existence asked about in \cite[Question~7]{CRnew}; it does not settle the
other alternative in that question.

\section{Formal verification and provenance}\label{sec:formal}
The companion development belongs to the \emph{InfinitaryCombinatorics}
library \cite{IC}. It reuses its definitions of AD, infinite MAD,
fin-sequences, retained traces, centeredness, and FI. Its new modules are in
\texttt{Formalizations/FIMAD/}. The toolchain remains Lean~4.30.0 with the
library's pinned mathlib revision. A pinned YesMetaZFC dependency supplies
the genuine first-order derivation and consistency notions \cite{YM}.

The checked combinatorial chain proves uniform trace coding, finite-error
trace obstruction, and both comparisons $\ap\leq\bnum\leq\anum$.
The comparison $\ap\leq\bnum$ is proved directly using binary-tree branches:
a weak separator for the eventually-zero branches and a family of infinite
branches would eventually bound the increasing enumerations of the latter.
The paper's exact cutoff and exclusion of all infinite FI MAD families
therefore have only the stated hypothesis $\s<\ap$; no additional cardinal
comparison is assumed. Arbitrary labels, including empty positive fibers,
are also covered. The development defines
$\ap$ by the least nonseparable AD cardinal and proves that its witness class
is nonempty; it does not replace $\ap$ by an unspecified parameter.

The complete construction in Section~\ref{sec:ch} is also checked.
Starting solely from $2^{\aleph_0}=\aleph_1$, it performs the well-founded
recursion on the actual first uncountable ordinal, proves maximality
against every infinite subset of $\omega$, and verifies FI against every
fin-sequence. The output has size $\aleph_1$. The positive extension,
finite ideal-cover, and finite-intersection preservation lemmas are
proved within the development, rather than supplied as assumptions.

There are substantive remaining formalization obligations. The BMZ forcing
construction, the relative consistency of CH, the stronger published positive
theorem under $\ap=\s=\cont$, and the topological results are not derived
in Lean here. Nor is the assertion
$\E$ encoded as a sentence of first-order set theory with its semantic
correspondence proved. The YesMetaZFC module checks only the general
proof-theoretic implication from consistent positive and negative
extensions to nonderivability on both sides. It is not an instantiation
of that implication for this problem. Accordingly, Theorem~\ref{thm:independence}
is a paper theorem using the named literature inputs, not a fully
machine-checked independence theorem. The source archive includes a
statement-by-statement coverage record and the verification procedure.

The research manuscript and this formalization were prepared with
generative AI assistance, including GPT-6 Astra. The author supplied the
research problem and coordinated the work. Checks performed by the
producing AI system do not constitute independent expert review; no such
review of this manuscript is claimed. The author is unaffiliated and
received no specific funding for this work. Known positive constructions,
cardinal comparisons, forcing results, and hyperspace theorems are credited
to their cited sources. No claim of an exhaustive priority search is made.

\appendix
\section{The same splitting number and continuum need not decide existence}
Corral and Rodrigues construct under CH a FI MAD family that survives
adding arbitrarily many Cohen reals; the countable-forcing case is
\cite[Theorem~5.3]{CR} and the arbitrary-Cohen extension is discussed in
Section~5 of that paper. This yields a useful comparison, separate from
the independence proof above.

Start from GCH with such a family $\A$ of size $\aleph_1$ and force with
$\operatorname{Add}(\omega,\omega_2)$. The forcing is ccc and preserves
cardinals. Counting nice names gives $\cont\leq\aleph_2$, and the distinct
coordinate reals give equality. A nice name for a real or a fin-sequence
has countable support, so it appears in a countable-coordinate extension.
The cited preservation theorem applies there. An infinite-index FI witness
continues to work in the full extension: the family has the same members,
and statements about finite intersections of its traces are absolute.
The same support argument checks maximality.

The first $\omega_1$ Cohen reals form a splitting family. For any infinite
$X$ in the final extension, choose a countable support for its name and
then an index below $\omega_1$ outside that support. That coordinate is
Cohen over the extension containing $X$, and the dense requirements to
put both bits at arbitrarily large positions in $X$ ensure splitting.
Thus $\s=\aleph_1$. Since the preserved MAD family still has size $\aleph_1$,
$\anum=\ap=\bnum=\aleph_1$ as well.
\[
\begin{array}{lcccc}
\text{Model} & \s & \ap & \cont & \E\\\hline
\text{Cohen extension above} & \aleph_1 & \aleph_1 & \aleph_2 & \text{true}\\
\text{BMZ model} & \aleph_1 & \aleph_2 & \aleph_2 & \text{false}
\end{array}
\]
In particular, the hypothesis $\ap=\cont$ cannot simply be dropped from
Corollary~\ref{cor:section}. This appendix relies on the cited preservation
result and is not included in the Lean verification.

\begin{thebibliography}{9}
\raggedright
\bibitem{CR} C\'esar Corral and Vinicius de O.~Rodrigues,
\emph{Fin-intersecting MAD families}, Filomat \textbf{38} (2024), no.~7,
2563--2578. \href{https://doi.org/10.2298/FIL2407563C}{doi:10.2298/FIL2407563C}.
Published numbering is used. Earlier version:
\href{https://arxiv.org/abs/2212.01484v2}{arXiv:2212.01484v2}.
\bibitem{BMZ} Taras Banakh, Micha\l{} Machura and Lubomyr Zdomskyy,
\emph{On critical cardinalities related to Q-sets}, Mathematical Bulletin
of the Shevchenko Scientific Society \textbf{11} (2014), 21--32.
\href{https://arxiv.org/abs/1306.0204v3}{arXiv:1306.0204v3}.
Theorem~2 is on p.~2 and Theorem~4(3) on p.~5 of that version.
\bibitem{CRnew} C\'esar Corral and Vinicius de O.~Rodrigues,
\emph{New examples of MAD families with pseudocompact hyperspaces},
Topology and its Applications \textbf{342} (2024), 108773.
\href{https://doi.org/10.1016/j.topol.2023.108773}{doi:10.1016/j.topol.2023.108773}.
Theorem~5 and Question~7 refer to
\href{https://arxiv.org/abs/2309.03405v2}{arXiv:2309.03405v2}.
\bibitem{IC} Haoxuan Ye, \emph{InfinitaryCombinatorics}, Lean source repository,
\url{https://github.com/EgoFakeFantasy/InfinitaryCombinatorics}.
The accompanying archive fixes the source and dependency hashes for this revision.
\bibitem{YM} \emph{YesMetaZFC}, Lean source repository,
\url{https://github.com/lanxinge/YesMetaZFC}, commit
\texttt{bae4fcc31b07b505986b11c6c2f13965ae6cd46d}.
\end{thebibliography}
\end{document}
